\chapter{Limitations and Future Work}\label{ch:limitations}

A study whose headline contribution is methodological honesty owes the reader an
equally honest account of its own boundaries. The findings reported in
\Cref{ch:experiments} are, for the most part, deliberately sober: at the scale we
study, a fairly tuned predictive-coding (PC) network is statistically
indistinguishable from a matched backpropagation (BP) baseline, and our systems
contribution---a fused CUDA settling kernel---wins only in a delimited batch regime.
None of these results should be over-read. This chapter therefore does two things.
First, it states precisely what our claims do \emph{not} entail: the scale at which
they hold, the seed budgets behind the confidence intervals, the architectural
regimes in which the kernel helps or hurts, and the exact provenance of the
``first-for-PC'' novelty claim. Second, it lays out the concrete extensions that
would either widen the scope of the present claims or test the hypotheses that this
study could only flank rather than resolve. The intent throughout is to leave a
reviewer with no surprises: every caveat that an adversarial reader might raise is
raised here first, by us.

\section{Limitations}

We group the limitations into four families: those of \emph{scientific scope} (what
the null result does and does not establish), those of \emph{the systems
contribution} (where the kernel wins and where it does not, and how far the novelty
claim reaches), those of \emph{the capability demonstrations} (the generative model,
the anomaly task, the incremental-PC variant), and those of \emph{experimental
reach} (scale and architecture). We treat each in turn.

\subsection{Scope of the scientific findings}

\subsubsection{The null result is scale- and budget-bounded, not universal}

The central empirical claim of this thesis is a \emph{null}: across bulk accuracy,
noise robustness, sample efficiency, domain-incremental learning, and
class-incremental learning, a loss-matched and per-method-tuned PC network shows no
advantage over BP that exceeds our measured noise (\Cref{sec:exp-pcbp}). It is
essential to state what this claim is \emph{not}. It is not a proof that PC can
never differ from BP. It is the considerably weaker---and we believe defensible---
statement that \emph{at MNIST and FashionMNIST scale, with a fairly tuned baseline
and a matched loss function, no PC advantage exceeds the seed-to-seed noise}. The
distinction matters because the literature reports advantages in regimes we did not
reach, and the theory \citep{innocenti2026limits} locates PC's distinctive behaviour
precisely in regimes (depth, width) that our $\lbrack 784, 256, 256, 10\rbrack$
network does not probe and that the theory itself predicts to become unstable as
they scale. Our result is a clean data point at one scale, not a refutation of the
field.

\subsubsection{The null does not globally refute Song et al.}

This caveat deserves its own treatment because it is the most likely point of
misreading. Our class-incremental experiments (\Cref{sec:exp-pcbp}) follow the
\emph{regime} studied by \citet{song2024prospective}, and \Cref{sec:exp-song}
reports an architecture-faithful replication of their \emph{exact} alternating
Split-FashionMNIST protocol: a four-layer $32$--$32$ sigmoid network with
Xavier-normal initialisation, two disjoint five-class tasks sharing one five-output
head, trained by alternating at the minibatch level. On the axes Song specifies
(depth, hidden width, activation, initialisation, task split, shared head,
alternation) we match the published method. Even so, two honest gaps remain. First,
the \emph{scale and seed budget} differ: we ran $n=10$ seeds and reported $68\%$
($1\sigma$) bootstrap intervals, whereas the original study operates at a different
scale with its own seed budget; a $1\sigma$ null at one scale is not a $95\%$ null
at another. Second, the original $84$-iteration training budget leaves every method
at chance in our hands---a faithfulness-versus-trainability tension that we resolved
by treating the training budget as a controlled axis from $84$ iterations to
convergence, rather than by silently enlarging it. Within that swept budget we find
no PC-versus-BP(MSE) difference exceeding $1\sigma$ at \emph{any} point, and raising
the seed count from five to ten did not promote the directional trend (PC slower
early, marginally ahead at convergence) to significance. The honest reading is
therefore: \emph{our most faithful replication agrees with our parity verdict rather
than overturning it}, while the residual differences in scale and budget mean we
cannot claim to have reproduced or refuted the original finding in its own terms.

\subsubsection{A corrected framing, not an original interpretation}

The interpretation we report in \Cref{sec:exp-song} is itself the product of an
adversarial three-lens review (literature faithfulness, statistics,
alternative-explanations) that caught an earlier draft misattributing a
learning-rate-robustness claim to \citet{song2024prospective}'s interference figure.
The original figure is a \emph{less-interference} claim with the learning rate
optimised per method, not an LR-robustness claim. We mention this here because it is
a limitation of method as much as a strength: it demonstrates that, absent the
adversarial check, we would have compared against the wrong baseline metric. The
safeguard is reusable and we recommend it, but its necessity is also a reminder that
single-pass readings of the literature in this field are error-prone, our own
included.

\subsubsection{The confound checklist explains, but does not exhaust}

We identify three confounds that manufacture spurious PC advantages in naive
comparisons (\Cref{sec:exp-pcbp}): the loss function (cross-entropy versus MSE, worth
$+5.9$ percentage points of bulk accuracy on its own), the
stability--plasticity trade-off in forgetting metrics, and an untuned baseline
learning rate. This list is offered as an explanation for the literature's
inconsistency, not as a closed taxonomy. Other confounds---initialisation schemes,
data ordering, the choice of optimiser for the baseline, the number of settling
steps $T$ relative to network depth---could equally distort a comparison, and we
controlled rather than exhaustively ablated several of them. The checklist is a
practical aid, not a completeness theorem.

\subsection{Limitations of the systems contribution}

\subsubsection{The kernel wins only in the small/medium-batch regime}

The fused settling kernel (\Cref{sec:kernel}, \Cref{sec:exp-kernel}) is honestly a
\emph{regime-specific} accelerator, not a general-purpose replacement for the
PyTorch backend. PC inference at MNIST scale is launch-overhead-bound: the PyTorch
backend issues $31$ kernel launches per settling step and runs at roughly
$0.5$\,ms/step independent of batch size, leaving the GPU $25$--$49\%$ utilised.
Fusing the entire $T$-step loop into a single launch removes that overhead and
yields $3.2\times$ at batch $64$. But the same property that makes the fused design
fast at small batch---holding the layer states resident on-chip across all $T$
steps---is architecturally opposed to the register-blocked tiling that a
cuBLAS-class GEMM needs: residency consumes precisely the register and shared-memory
budget a fast matrix multiply requires. Consequently, batch-tiling and
input-activation caching push the crossover only to roughly $1024$--$1500$, and at
batch $\geq 2048$, where \emph{compute} rather than launch overhead dominates, cuBLAS
wins and the correct backend is \texttt{pytorch}. We report this not as ``not yet
optimised'' but as a genuine systems finding---the fused-resident and tiled-GEMM
designs are in tension---yet it is unambiguously a limitation on the kernel's
applicability.

\subsubsection{The large-batch tiled path is still depth-two}

The kernel now exists in two forms with different generality. The per-sample path
(one thread block per sample) has been generalised to \emph{arbitrary depth}: the
per-layer weights, states, and biases are passed as device-pointer arrays with
shared-memory offsets computed at launch, and the result is validated against the
PyTorch equilibrium to roughly $10^{-6}$ across depths two through five (one to four
hidden layers), for tanh and sigmoid activations and batches of $32$ and $256$. The
batch-tiled path (the one that buys large-batch weight reuse), however, remains
specialised to the two-hidden-layer topology. The general-depth kernel is therefore
restricted to the small-batch, launch-bound regime---which is the regime in which it
wins, so the restriction is not crippling---but a register-blocked tiled GEMM for the
large-batch \emph{and} deep regime is not implemented. Lifting the depth
specialisation on the tiled path is reserved future work.

\subsubsection{``First for PC'' carries the to-our-knowledge caveat}

The novelty claim is stated carefully and we restate its exact scope here. To our
knowledge, this is the first hand-written, fused, single-launch CUDA kernel that runs
the entire $T$-step PC settling loop with layer states held resident on-chip across
steps. This is supported by a direct source audit of eleven open PC libraries (PCX,
JPC, ngc-learn, $\mu$PC, PRECO, Torch2PC, pypc, pybrid, ePC, and Song's own
repositories), none of which ship a custom CUDA settling kernel---they are uniformly
JAX/JIT or PyTorch-autograd---and by an additional audit of the
Equilibrium-Propagation family, the closest algorithmic relative, which likewise
ships no hand-written CUDA settling kernel (the one EP repository with hand-written
GPU kernels uses per-step elementwise Triton kernels with the settling iteration on
the host). Two qualifications are mandatory. First, the \emph{technique}---a fused,
single-launch, states-resident persistent kernel---is established prior art, dating
to cuDNN's persistent RNN and Baidu's \texttt{persistent-rnn}
\citep{diamos2016persistent} and re-applied in recent LLM megakernels; our
contribution is its transfer to, and measurement on, the PC settling loop, not the
invention of the pattern. The closest fixed-point analogue, deep equilibrium models
\citep{bai2019deq}, and the closest PC-side work, error-optimisation PC
\citep{goemaere2025epc}, attack related goals without a custom kernel and are cited
and delimited accordingly. Second, a negative existence claim over unindexed or
industrial code cannot be proved; the ``to our knowledge'' caveat is therefore not
a rhetorical hedge but a literal statement of the evidence's limit. The observation
that PC inference is launch-overhead-bound is itself a generic ``framework tax'' and
is \emph{not} claimed as a contribution.

\subsection{Limitations of the capability demonstrations}

\subsubsection{Generative output is class prototypes, not diverse samples}

The generative PC of \Cref{sec:exp-gen} genuinely recovers the ``one mechanism, many
tasks'' property: the same settling rule classifies, generates, inpaints, and detects
anomalies, depending only on what is clamped. But the generative model maps a
one-hot label to an image, so its ten possible inputs produce at most ten distinct
outputs: recognisable \emph{per-class prototypes} (effectively the per-class mean),
not a diverse sample distribution. Inpainting is correspondingly prototype-driven and
blurry rather than pixel-sharp. Producing diverse samples would require a continuous
latent with a prior (a VAE-like construction), which we did not build. The
demonstration establishes the qualitative capability, not generative fidelity.

\subsubsection{Uniform noise is an easy out-of-distribution test}

The anomaly-detection result---an AUC of $1.0$ separating in-distribution MNIST
(residual energy $5.1$) from uniform-noise inputs (energy $94.5$)---is real but must
not be over-read. Uniform noise is an \emph{easy} out-of-distribution (OOD) signal:
it lies far from any plausible image manifold, so perfect separation is a weak test
of the detector. A genuinely informative evaluation would use \emph{near}-OOD inputs,
such as FashionMNIST presented to an MNIST-trained model, or rotated digits, where
the OOD distribution shares low-level statistics with the training data. We report
the easy result honestly and flag the harder test as the meaningful one.

\subsubsection{Incremental PC is learning-rate sensitive}

The incremental-PC (iPC) variant \citep{salvatori2024ipc} of \Cref{sec:exp-pcbp}
interleaves a weight update with every settling step rather than updating once at
equilibrium. Because it applies $T$ weight updates per minibatch, its effective
weight learning rate is roughly $T$ times larger, and at the standard rate it
diverges. At a rate scaled down by approximately $1/T$ it is stable and, in a
fast-training regime, outperforms standard PC ($74.3\%$ versus $63.1\%$ on a short
MNIST run). The capability is therefore real but conditional: iPC trades robustness
to the learning rate for faster per-step learning, and a practitioner who omits the
$\sim\!\eta_w/T$ rescaling will simply see it diverge. We report it as an optional
extension with this sensitivity made explicit, not as a drop-in improvement.

\subsubsection{Precision schedules remain reserved}

Throughout this work the precision (inverse-variance) weighting $\Pi$ of the
prediction errors is fixed to the identity, which reduces the energy
$E = \tfrac12 \sum_k \lVert \varepsilon_k \rVert^2$ to a plain sum of squared
errors---the modern default of PCX and JPC \citep{pinchetti2024pcx, millidge2024jpc}.
Learnable or scheduled precisions are a documented dimension of the PC objective and
a plausible lever on both the depth-scaling problem (precision-weighted gradients are
one of the proposed fixes) and the energy imbalance across layers, but we did not
implement them. The ``precision schedule'' search dimension is therefore reserved,
and any claim we make is implicitly conditioned on $\Pi = I$.

\subsection{Limitations of experimental reach}

\subsubsection{Scale and the gap to Transformers}

The entire study lives at MNIST and FashionMNIST scale with a vanilla
state-optimisation PC network. We do not demonstrate PC on larger datasets, on
convolutional or attention architectures, or at anything approaching the scale at
which the field's open questions (\Cref{ch:related}) actually bite. The
\emph{scale gap}---PC has not been shown to work at Transformer scale---is untouched
here and is frontier-laboratory work; our contribution is a clean, well-characterised
building block, not a scaling result.

\subsubsection{Depth scaling is sidestepped, not solved}

Closely related, the depth-scaling failure of PC---the observation that deeper
PC-trained networks are often \emph{worse}, the opposite of BP, driven by an
exponential signal-decay problem \citep{goemaere2025epc}---is \emph{characterised but
not solved} here. We reproduce it cleanly (\Cref{sec:exp-decay}): standard SO-PC matches
BP at one--two hidden layers but collapses to near-chance by eight, with the per-layer
learning signal decaying $\sim\!669\times$ from output to input. We also show that a
strictly local precision schedule $\Pi_k=\gamma^{(n-k)}$ mechanistically mitigates it
(flattening the signal profile $10$--$30\times$ and recovering much of the accuracy at
moderate depth), but only partially at depth eight, and---crucially---this is a
\emph{known} technique \citep{qi2025settling, innocenti2025mupc}, not our contribution,
and it is bettered by those works and by error optimisation \citep{goemaere2025epc}.
What remains genuinely open for us is full deep-PC \emph{trainability}: matching BP at
the depths where the published fixes operate, which is frontier-laboratory work. The
problem is now mapped at our scale, not resolved.

\subsubsection{Single optimiser regime and offline evaluation}

Finally, the comparisons use a single optimisation regime (plain local updates for PC,
SGD-style updates for the matched BP baseline) and an offline, batch-trained
evaluation. We do not study the strict online (batch-size-one) regime that is among
PC's most cited putative niches, nor optimiser interactions (momentum, adaptive
methods) that could shift the parity verdict in either direction. All experiments are
implemented in pure PyTorch with hand-derived updates verified against autograd as an
external oracle, and the full suite comprises forty offline tests; this gives
correctness confidence but is not a substitute for the broader sweep just described.

\section{Future Work}

The limitations above map naturally onto a future-work programme. We organise it by
the same two threads as the thesis---the systems contribution and the scientific
contribution---and add the capability extensions that the generative results invite.
Where a direction is squarely frontier-laboratory work we say so; where a solo,
small-scale project can make a real dent we say that too.

\subsection{Extending the systems contribution}

\subsubsection{A register-blocked tiled GEMM and an ahead-of-time build}

The clearest engineering next step is to attack the large-batch regime directly.
The kernel currently loses to cuBLAS at batch $\geq 2048$ because its in-kernel
matrix multiply uses scalar fused-multiply-adds with no register blocking and reaches
only a fraction of the FLOP peak. A register-blocked, shared-memory-tiled GEMM---in
effect reimplementing the core of cuBLAS inside the fused kernel---would, if
combined with a careful budgeting of the residency-versus-tiling trade-off described
above, narrow or close the large-batch gap. This is a substantial undertaking and
arguably of questionable return given that cuBLAS already serves that regime well,
but it is the principled way to lift the fused design beyond its current ceiling. In
parallel, replacing the current just-in-time \texttt{torch.utils.cpp\_extension.load}
build with an ahead-of-time (AOT) \texttt{setuptools}/\texttt{CUDAExtension} build
against the ABI-stable LibTorch API would improve reproducibility and packaging.
Lifting the batch-tiled path to arbitrary depth (it remains depth-two) is the
companion task.

\subsubsection{An Nsight Compute profile to corroborate the launch-overhead account}

Our mechanistic evidence that launch overhead is the addressed bottleneck rests on a
profiler launch count (the PyTorch backend issues $620$, $1240$, or $2480$ launches
for $T = 20, 40, 80$ versus the fused kernel's single launch). A full Nsight Compute
report---occupancy, achieved bandwidth, memory-versus-compute roofline placement,
warp-stall reasons---would harden this account against the ``solver-view'' objection
raised by ODE-solver-based libraries such as JPC and ePC, demonstrating at the
hardware-counter level that the PyTorch path is launch-bound and the fused path is
not. We note the launch-count figure is already presented in \Cref{sec:exp-kernel};
the Nsight report is the natural deepening.

\subsection{Extending the scientific contribution}

\subsubsection{An Equilibrium-Propagation comparison arm}

Equilibrium Propagation (EP) \citep{scellier2017ep, laborieux2021ep} is the closest
algorithmic relative of PC: both relax to an equilibrium and derive a local weight
update from it. We audited EP for the novelty claim but did not implement it as a
comparison arm. Adding an EP arm under the same fair-comparison protocol---matched
architecture, initialisation, data, and loss; per-method tuning; bootstrap
confidence intervals---would situate our PC-versus-BP parity verdict within the
broader energy-based-learning family and test whether EP, too, collapses to BP-like
behaviour at this scale or carves out a distinct regime.

\subsubsection{Probing where PC and BP diverge}

The deepest open question we can only flank is the PC$\,\approx\,$BP equivalence
puzzle: the theory \citep{innocenti2026limits} holds that the width- and
depth-stable feature-learning parametrisations of PC coincide with those of BP, and
that PC's distinctive regime becomes unstable as it scales---leaving open what PC is
ultimately \emph{for}. A small project cannot resolve this, but it can contribute
data points. Concretely, one could test whether the PC and BP \emph{solutions}
themselves diverge---do they reach different minima, exhibit different loss-landscape
curvature \citep{zahid2023curvature}, or differ in robustness---even when their test
accuracies match. A measurable divergence in solution geometry, despite accuracy
parity, would be a genuine contribution to the question, and our matched-init harness
(the BP baseline clones the PC initialisation) is well suited to it.

\subsubsection{Harder near-OOD and continual-learning baselines}

On the empirical side, the anomaly task should move from easy uniform-noise OOD to
near-OOD (FashionMNIST against an MNIST-trained model, rotated digits), which would
turn the AUC $1.0$ result into an informative measurement. On the
continual-learning side, our EWC baseline \citep{kirkpatrick2017ewc} already behaves
as a positive control---it forgets less than the matched BP, with backward transfer
of $-41.1\%$ versus $-45.7\%$---and the natural next step is to add a replay-based
baseline such as gradient episodic memory \citep{lopezpaz2017gem} and a joint-training
upper bound, ideally standardised through a continual-learning framework such as
Avalanche \citep{lomonaco2021avalanche}, to locate the absolute forgetting numbers
rather than only comparing PC against BP.

\subsection{Extending the capability demonstrations}

\subsubsection{A diverse generative latent}

To move the generative PC from class prototypes to diverse samples, the one-hot label
input must be replaced by a continuous latent with a prior, settled as a free
variable under a generative training objective in the spirit of the original
predictive-coding model \citep{rao1999predictive}. This is a distinct
model variant rather than a tuning change, and it is the most direct route to a
generative PC that produces sample \emph{diversity} and supports the sharper
inpainting that prototype-driven reconstruction cannot.

\subsubsection{Depth-scaling diagnostics and proposed fixes at small scale}

Finally, on the depth-scaling problem we have already taken the small-project step of
\emph{characterising} it precisely (\Cref{sec:exp-decay}): the ``deeper-is-worse''
effect is reproduced across two, four, six, and eight hidden layers, the per-layer
learning-signal decay is measured and plotted, and one local fix (a precision schedule)
is shown to mechanistically flatten the imbalance --- though it is a known technique and
only partial at depth. The natural continuation is to implement error optimisation
\citep{goemaere2025epc} as a fourth arm and complete the benchmark matrix (PyTorch-SO /
fused-SO / EO / BP, over accuracy, wall-clock, and per-layer signal versus depth), which
would turn the present ablation into a full comparative study of \emph{how} each method
keeps deep PC trainable; reaching backpropagation-equivalent depth in a local formulation
remains the frontier piece beyond it.

\medskip

Taken together, these directions do not change the thesis's sober verdict; they
sharpen its edges. The systems extensions would widen the kernel's regime of
usefulness without retracting the architectural trade-off we identified; the
scientific extensions would test the parity verdict against a wider family of methods
and probe the one question---where, if anywhere, PC genuinely diverges from BP---that
this study could only frame, not answer. That this list is long is itself in keeping
with the project's stance: an honest baseline is valuable precisely because it makes
the next questions clear.
